%  =============================================================================
%  Datei-Hinweise & Konfiguration: siehe config.md
%  -> Abschnitt "anhang/ki-prompts/tool-fable.tex"
%  =============================================================================

\section{Anthropic Claude (Fable~5)}
\label{sec:prompts-fable}

\kiwerkzeug{Anthropic Claude, Modell Fable 5}{Modellbezeichner
\texttt{claude-fable-5}, veröffentlicht am 09.06.2026}{https://claude.ai}
% Eintrag im KI-Verzeichnis (Leitfaden Kap. 6.2.4) erzwingen mittels \nocite.
\nocite{claudeFable2026}

Eingesetzt zur Unterstützung bei der Implementierung: Python- und
Shell-Skripte des Agenten, die Cloud-Run-Anwendung des öffentlichen
Dashboards, die Gegenprüfung der zuvor erzeugten Testskripte sowie die beiden
als KI-generiert gekennzeichneten Abbildungen dieser Arbeit
(\cref{fig:zielarchitektur} und \cref{fig:projektstruktur}).

\subsection{Unterstützung bei Python- und Shell-Skripten}
\label{sec:prompts-fable-skripte}
\begin{enumerate}
    \item \enquote{Überarbeite die folgende Funktion so, dass ein
            Verbindungsabbruch des Datenstroms nicht zum Prozessabbruch führt,
            sondern erkannt, protokolliert und über einen Wächterprozess
            neu verbunden wird; erhalte dabei die bestehende Signatur:
            [\dots]}
            \par\hfill\textit{Kontext: \cref{sec:umsetzung-architektur}}

    \item \enquote{Schreibe ein Shell-Skript, das den Entscheidungszyklus
            startet, eine Prozesssperre je Konto setzt, überlappende Läufe
            verhindert und im Fehlerfall mit eindeutigem Rückgabewert endet:
            [\dots]}
            \par\hfill\textit{Kontext: \cref{sec:umsetzung-implementierung}}

    \item \enquote{Erkläre, welche Fehlerbilder verteilter Systeme der
            nachfolgende Ausführungspfad noch nicht abfängt, und schlage je
            Fall die konservativste Behandlung vor: [\dots]}
            \par\hfill\textit{Kontext: \cref{sec:umsetzung-implementierung}}
\end{enumerate}

\subsection{Öffentliches Dashboard (Cloud-Run-Anwendung)}
\label{sec:prompts-fable-dashboard}
\begin{enumerate}
    \item \enquote{Entwirf einen zustandslosen Dienst, der redigierte
            JSON-Feeds aus einem Objektspeicher liest und daraus eine
            Weboberfläche rendert, ohne dass die handelnden Maschinen einen
            eingehenden Port öffnen müssen: [\dots]}
            \par\hfill\textit{Kontext: \cref{sec:umsetzung-implementierung}}

    \item \enquote{Erzeuge die Oberfläche des Dashboards so, dass der
            Indexverlauf mit den Durchschnittslinien und den Regimewechseln
            dargestellt wird und je Entscheidung das aggregierte Votum mit
            aufklappbaren Einzelbegründungen erscheint: [\dots]}
            \par\hfill\textit{Kontext: \cref{sec:umsetzung-implementierung}}

    \item \enquote{Prüfe den folgenden Rendering-Pfad darauf, ob Inhalte aus
            dem Journal ungeprüft in das Dokument geschrieben werden, und
            schlage eine Ausgabekodierung vor: [\dots]}
            \par\hfill\textit{Kontext: \cref{sec:umsetzung-implementierung}}
\end{enumerate}

\subsection{Verifikation der erzeugten Testskripte}
\label{sec:prompts-fable-testverifikation}

Die Testskripte wurden mit Claude Opus~5 entworfen
(\cref{sec:prompts-claude-tests}); die nachfolgenden Prompts dienten allein
ihrer unabhängigen Gegenprüfung durch ein zweites Modell.

\begin{enumerate}
    \item \enquote{Prüfe die folgende Testsuite kritisch: Welche Verzweigung
            des Produktivcodes wird von keinem Testfall erreicht, und welcher
            Test würde auch dann bestehen, wenn die geprüfte Zusicherung
            entfiele? [\dots]}
            \par\hfill\textit{Kontext: \cref{sec:umsetzung-implementierung}}

    \item \enquote{Benenne Testfälle, die nur scheinbar prüfen, was ihr Name
            behauptet, und schlage jeweils eine Formulierung vor, die den
            tatsächlich geprüften Sachverhalt trifft: [\dots]}
            \par\hfill\textit{Kontext: Anhang~\ref{sec:anhang-archiv-tests}}
\end{enumerate}

\subsection{Erzeugung der KI-generierten Abbildungen}
\label{sec:prompts-fable-abbildungen}

Beide Abbildungen sind in der jeweiligen Bildunterschrift als KI-generiert
ausgewiesen (Leitfaden Kap.~5.7 in Verbindung mit Kap.~6.2.4). Inhaltliche
Grundlage waren ausschließlich der eigene Quellcodebestand und die eigene
Architekturplanung; die Modellantwort lieferte die grafische Umsetzung.

\begin{enumerate}
    \item \enquote{Erzeuge eine SVG-Grafik der Zielarchitektur des
            Handelsagenten auf Basis der bereits erfolgten Implementierungen
            sowie der geplanten Bestandteile; umgesetzte Elemente
            durchgezogen, geplante gestrichelt, mit Legende. Style die Grafik
            attributbasiert ohne \texttt{<style>}-Block, damit die
            Inkscape-Konvertierung verlustfrei arbeitet: [\dots]}
            \par\hfill\textit{Kontext: \cref{fig:zielarchitektur}}

    \item \enquote{Erzeuge aus der nachfolgenden ASCII-Repräsentation der
            Verzeichnis- und Dateistruktur eine SVG-Grafik, die die Struktur
            in einer monospaced Schriftart schwarz auf weiß rendert, die
            Ebenen-Header durch Fettsatz hervorhebt und je Datei die
            Herkunftsmaschine kennzeichnet: [\dots]}
            \par\hfill\textit{Kontext: \cref{fig:projektstruktur}}
\end{enumerate}

%  =============================================================================
%  Ende
%  =============================================================================
